\section{What The Strong Evidence Actually Buys}

The 2024 Communications Biology source matters because it supports the claim
that conscious states come with organized dynamical signatures. For this
manuscript, that is mainly a governance result: the book is allowed to speak
about structured state organization without sounding biologically naive.

The 2025 Communications Biology source adds a different layer of value. By
connecting cortical connectivity, local dynamics, and stability correlates of
global conscious states, it supports the idea that persistence depends on both
coupling architecture and local system behaviour. In business terms, this is
the bridge from everyday ``current mood'' talk to a more serious ``stability
architecture'' vocabulary.

The Nature Neuroscience sleep-transition source is the most useful transition
anchor because it deals with a concrete biological regime shift rather than a
moral story about effort. That makes bifurcation-style language relevant to the
book's later discussion of decision windows and leverage points.

\begin{discussion}
At manuscript level, the strong evidence buys four permissions:
\begin{itemize}[nosep]
  \item permission to talk about organized state structure rather than loose
        introspective mood labels;
  \item permission to treat persistence as an architectural problem rather than
        an accident of personality;
  \item permission to treat transitions as potentially sharp,
        parameter-sensitive events;
  \item permission to discuss leverage in terms of regime vulnerability and
        stability loss.
\end{itemize}
\end{discussion}

\begin{remark}[What the evidence does \emph{not} buy]
The same sources do not authorize claims that neuroscience has directly
measured free will, that every behavioural frame has a known neural
realization, or that simply borrowing a word such as ``bifurcation'' proves a
behavioural model. The chapter gains credibility only if it refuses proof by
vocabulary.
\end{remark}
